\section{Atividades Paralelas Desenvolvidas pelo Bolsista}

Durante o período correspondente aos dois primeiros trimestres de execução
do projeto, o bolsista participou de atividades acadêmicas complementares
voltadas à consolidação da formação científica e técnica em computação
quântica, contribuindo para o aprofundamento dos conhecimentos necessários
ao desenvolvimento da pesquisa. A documentação comprobatória das atividades
descritas nesta seção encontra-se disponível em material suplementar online
\cite{leao2026material_suplementar}.

\subsection{Participação em Evento Científico}

O bolsista participou, na condição de ouvinte, do II Simpósio de Computação
e Informação Quântica, promovido pela Liga Acadêmica de Computação e
Informação Quântica (LACIQ), realizado na Universidade Federal de Pernambuco
(UFPE), no período de 13 a 21 de novembro de 2025. O evento integrou as
comemorações dos 100 anos da Primeira Revolução Quântica e contou com
palestras nacionais e internacionais abordando temas relacionados à
computação quântica, informação quântica e química computacional quântica.

A participação proporcionou contato direto com pesquisadores atuantes na
área, ampliando a compreensão do estado atual das pesquisas e aplicações
científicas associadas ao tema do presente projeto. A carga horária total
do evento foi de 20 horas, conforme certificados disponíveis no material
suplementar \cite{leao2026material_suplementar}, cuja autenticidade pode
ser verificada por meio do sistema Even3.

\subsection{Cursos e Certificações em Computação Quântica}

Como atividade complementar de capacitação técnica, o bolsista realizou
cursos de formação oferecidos pela plataforma IBM Quantum Learning,
voltados aos fundamentos teóricos e práticos da computação quântica e
algoritmos variacionais.

Entre as certificações obtidas, destacam-se:

\begin{itemize}
\item Basics of Quantum Information;
\item Fundamentals of Quantum Algorithms;
\item Quantum Computing in Practice;
\item Quantum Chemistry with VQE;
\item Variational Algorithm Design.
\end{itemize}

Os cursos abordam conceitos fundamentais como representação da informação
quântica, construção de algoritmos quânticos, execução prática em hardware
quântico e implementação de algoritmos variacionais aplicados à química
quântica, diretamente relacionados às atividades desenvolvidas no presente
projeto. Os certificados correspondentes encontram-se disponibilizados
no material suplementar online \cite{leao2026material_suplementar}.